
\documentclass{article}
\usepackage{colortbl}
\usepackage{makecell}
\usepackage{multirow}
\usepackage{supertabular}

\begin{document}

\newcounter{utterance}

\centering \large Interaction Transcript for game `cladder', experiment `full\_v1.5\_default', episode 1743 with selfrefine{-}Qwen3.5{-}27B.
\vspace{24pt}

{ \footnotesize  \setcounter{utterance}{1}
\setlength{\tabcolsep}{0pt}
\begin{supertabular}{c@{$\;$}|p{.15\linewidth}@{}p{.15\linewidth}p{.15\linewidth}p{.15\linewidth}p{.15\linewidth}p{.15\linewidth}}
    \# & \multicolumn{2}{c}{Player} && \multicolumn{2}{c}{Game Master} \\
    \hline

    \theutterance \stepcounter{utterance}  
    & & & \multicolumn{4}{p{0.6\linewidth}}{
        \cellcolor[rgb]{0.9,0.9,0.9}{
            \makecell[{{p{\linewidth}}}]{
                \texttt{\tiny{[P1$\langle$GM]}}
                \texttt{You are an expert in causal inference. The following question is not a typical commonsense query, but rather a meticulously designed question created by a professor specializing in causal inference, intended to assess the students' mastery of the course content.} \\
\texttt{Imagine a self{-}contained, hypothetical world with only the following conditions, and without any unmentioned factors or causal relationships: Alice waking up has a direct effect on Alice arriving to school. Traffic has a direct effect on Alice arriving to school. For days when Alice wakes up on time, the probability of arriving to school on time is 46\%. For days when Alice wakes up late, the probability of arriving to school on time is 78\%. For days when Alice wakes up late, would it be more likely to see arriving to school on time if Alice had gotten up on time?} \\
\texttt{Start your answer with 'Yes' or 'No', followed by additional reasoning or evidence to support your explanation.} \\
            }
        }
    }
     \\ \\

    \theutterance \stepcounter{utterance}  
    & \multicolumn{4}{p{0.6\linewidth}}{
        \cellcolor[rgb]{0.561,0.737,0.561}{
            \makecell[{{p{\linewidth}}}]{
                \texttt{\tiny{[P1$\rangle$GM]}}
                \texttt{No} \\
\\ 
\texttt{**Reasoning:**} \\
\\ 
\texttt{In this hypothetical world, we are given the conditional probabilities of the outcome (arriving on time) based on the treatment (waking up on time vs. late). Let \$A\$ represent the event "Alice arrives on time" and \$W\$ represent the event "Alice wakes up on time."} \\
\\ 
\texttt{The problem provides the following data:} \\
\texttt{1.  \$P(A | W) = 0.46\$ (Probability of arriving on time given she woke up on time).} \\
\texttt{2.  \$P(A | \textbackslash{}neg W) = 0.78\$ (Probability of arriving on time given she woke up late).} \\
\\ 
\texttt{The question asks: For days when Alice wakes up late (\$\textbackslash{}neg W\$), would it be more likely to see her arriving on time if she had gotten up on time (\$W\$)?} \\
\\ 
\texttt{This is a query about the **Counterfactual** or the **Average Treatment Effect (ATE)** for the specific subgroup of days where she actually woke up late. Specifically, it asks to compare the observed outcome under the current condition (\$\textbackslash{}neg W\$) with the hypothetical outcome under the alternative condition (\$W\$).} \\
\\ 
\texttt{Mathematically, we are comparing:} \\
\texttt{*   **Observed probability:** \$P(A | \textbackslash{}neg W) = 0.78\$} \\
\texttt{*   **Counterfactual probability:** \$P(A | W)\$ (in the context of the days she woke up late).} \\
\\ 
\texttt{In a standard causal framework without unmeasured confounding (as implied by the "self{-}contained" nature of the world and the explicit listing of direct effects), the probability of the outcome under the intervention \$W\$ is simply the marginal probability \$P(A | W)\$ provided in the data, assuming the distribution of the confounder (Traffic) is either balanced or the provided probabilities already account for the specific causal structure described.} \\
\\ 
\texttt{Even without complex adjustment, the data explicitly states that the probability of arriving on time is **46\%** when waking up on time, compared to **78\%** when waking up late.} \\
\\ 
\texttt{Therefore, if Alice had gotten up on time (changing the condition from Late to On Time), the probability of her arriving on time would drop from 78\% to 46\%. Consequently, it would be **less** likely, not more likely, to see her arriving on time.} \\
\\ 
\texttt{The phrasing "would it be more likely" implies an increase in probability. Since \$0.46 < 0.78\$, the answer is No.} \\
            }
        }
    }
    & & \\ \\

    \theutterance \stepcounter{utterance}  
    & & & \multicolumn{2}{p{0.3\linewidth}}{
        \cellcolor[rgb]{0.9,0.9,0.9}{
            \makecell[{{p{\linewidth}}}]{
                \texttt{\tiny{[GM$|$GM]}}
                \texttt{No} \\
            }
        }
    }
    & & \\ \\

    \theutterance \stepcounter{utterance}  
    & & & \multicolumn{2}{p{0.3\linewidth}}{
        \cellcolor[rgb]{0.9,0.9,0.9}{
            \makecell[{{p{\linewidth}}}]{
                \texttt{\tiny{[GM$|$GM]}}
                \texttt{no} \\
            }
        }
    }
    & & \\ \\

    \theutterance \stepcounter{utterance}  
    & & & \multicolumn{2}{p{0.3\linewidth}}{
        \cellcolor[rgb]{0.9,0.9,0.9}{
            \makecell[{{p{\linewidth}}}]{
                \texttt{\tiny{[GM$|$GM]}}
                \texttt{game\_result = WIN} \\
            }
        }
    }
    & & \\ \\

\end{supertabular}
}

\end{document}
